\chapter{Inflation, Structure, and the Loss of Initial Conditions}

\section{Omega and the Evidence for Flatness}

A student's question about big \(\Omega\) sends the discussion back to the Friedmann-Robertson-Walker equation. For the present era, well after radiation domination, the expansion law is written with a vacuum-energy term and a matter term, while radiation is dropped because it is now very small:
\begin{equation}
H^2=\left(\frac{\dot a}{a}\right)^2=\frac{8\pi G}{3}\left(\rho_0+\rho_m\right)-\frac{k}{a^2}.
\end{equation}
The vacuum term does not dilute. The matter term does:
\begin{equation}
\rho_m=\frac{M}{a^3}.
\end{equation}
\(\Omega\) is the fraction of the expansion equation carried by the different components.

Divide by \(H^2\). Then the normalized terms add to one. In a regularized notation,
\begin{equation}
\Omega_0=\frac{8\pi G\rho_0}{3H^2},\qquad
\Omega_m=\frac{8\pi G\rho_m}{3H^2},\qquad
\Omega_k=-\frac{k}{a^2H^2},
\end{equation}
and therefore
\begin{equation}
\Omega_0+\Omega_m+\Omega_k=1.
\end{equation}
\editorialnote{The audio and blackboard are rough in this stretch, and the speaker shifts among several names for the vacuum contribution. The notation here regularizes the argument after division by \(H^2\).}
If space is flat, \(k=0\), so
\begin{equation}
\Omega_0+\Omega_m=1.
\end{equation}
If the matter-plus-vacuum sum is greater than one, the curvature is positive. If it is less than one, the curvature is negative.

The observational logic is then straightforward. Measure ordinary matter and dark matter. Measure the Hubble expansion. Measure the curvature geometrically on very large scales. Notice also that the universe is accelerating. The sign of the vacuum term is not \(H\) by itself but the tendency of \(\dot a/a\) toward a constant, which means that the scale factor approaches the exponential branch,
\begin{equation}
a(t)\propto e^{Ht}.
\end{equation}
The several measurements cross-check one another and say that the universe is very nearly flat, to within a few percent.

\begin{classroomqa}
\lecturetimestamp{00:10:50}
\audiencequestion{Does matter include dark matter?}
\lecturerresponse{Yes. For the present discussion the matter term includes dark matter. One could also add radiation, but today that fraction is so small that it can be ignored.}
\end{classroomqa}

\begin{classroomqa}
\lecturetimestamp{00:11:11}
\audiencequestion{Is there some measure of dark energy other than the Hubble constant?}
\lecturerresponse{The point is not the value of \(H\) by itself, but the way the expansion tends toward an exponential law. Together with the matter measurement and the geometric measurement of flatness, that lets one infer the vacuum term.}
\end{classroomqa}

\section{From Flatness to Structure}

Why is the universe so flat? The proposal is inflation. Flat means very big, very smooth, and very homogeneous on scales of billions, really tens of billions, of light years. Exponential expansion can stretch a region until it looks that way.

But one can overdo a good thing. A perfectly smooth universe stays perfectly smooth. Gravity does not manufacture lumpiness out of exact symmetry. If structure exists now, then structure had to begin from a small inhomogeneity earlier. Tracing the present universe backward, the inhomogeneity at the time of transparency comes out at roughly one part in \(10^5\):
\begin{equation}
\frac{\delta\rho}{\rho}\sim 10^{-5}.
\end{equation}
The universe was extremely smooth, but not perfectly so.

From there gravity does what ordinary forces in a room full of air do not do. A small overdensity pulls matter toward itself. The nearby region is depleted, so the underdensity becomes more underdense and the overdensity becomes more overdense. Gravity reinforces the contrast instead of erasing it.

The natural question is why this did not happen earlier. The answer is pressure. Earlier, radiation pressure was large. Using
\begin{equation}
p=w\rho,
\end{equation}
the familiar cases are
\begin{equation}
w_{\rm rad}=\frac13,\qquad
w_m=0,\qquad
w_{\rm vac}=-1.
\end{equation}
Positive pressure homogenizes. It pushes disturbances out and smooths the medium. Matter, treated on large scales as pressureless, offers no such support. With pressure present, gravity is held off. When pressure becomes negligible, gravity wins.

Dark matter is crucial in the order of events. The first collapsing lumps are dark-matter halos. Ordinary baryonic matter then falls into those halos and, because baryons can dissipate, collects nearer the center. The dark matter does not dissipate in the same way, so it remains spread through the halo. The beginning of that contraction is slightly earlier than transparency, by roughly a factor of ten in time, which is why the microwave background does not show the onset directly.

A brief aside fixes the later chemical history. Almost all elements heavier than helium are products of stellar evolution rather than early-universe cosmology. The early universe makes essentially hydrogen, helium, and small traces of a few light nuclei.

What is observed on the sky is not mass density itself but temperature anisotropy in the microwave background. Density fluctuations contribute to it. Velocity fluctuations contribute too. Overdense regions tend to look slightly colder because photons lose extra energy climbing out of deeper gravitational wells, while Doppler shifts from the motion of the gas are mixed in on top of that. The full numerical analysis is complicated. The robust statement is simpler: by the time the universe became transparent, the energy density varied by about one part in \(10^5\).

\begin{classroomqa}
\lecturetimestamp{00:20:29}
\audiencequestion{Why didn't gravity make the lumps grow earlier?}
\lecturerresponse{Because earlier the universe was radiation dominated, and radiation carried enough positive pressure to homogenize the medium. Only when radiation became unimportant compared with matter did gravity begin to dominate the growth of the lumps.}
\end{classroomqa}

\begin{classroomqa}
\lecturetimestamp{00:24:35}
\audiencequestion{Don't you need a dissipation mechanism for collapse?}
\lecturerresponse{Ordinary baryonic matter does need dissipation if it is to settle toward the center. But dark matter can begin by making halos without that sort of collisional dissipation. There is also a cosmological kind of damping tied to expansion itself, and that will matter later.}
\end{classroomqa}

\begin{classroomqa}
\lecturetimestamp{00:30:34}
\audiencequestion{Were there also velocity fluctuations, so that later galaxies could acquire angular momentum?}
\lecturerresponse{Yes. The primordial pattern involved density fluctuations and velocity fluctuations. What is seen on the sky is a temperature pattern produced by both the density variations and the motion of the gas.}
\end{classroomqa}

\section{The Inflaton and the Plateau}

The mechanism now introduced is a scalar field \(\phi\). Its energy density has three pieces:
\begin{equation}
\rho_\phi=\frac12\dot\phi^2+\frac12(\nabla\phi)^2+V(\phi).
\end{equation}
Time variation contributes. Spatial variation contributes. A scalar field can also carry a true potential energy \(V(\phi)\). That part is not exotic; the Higgs field is the familiar example. What matters here is the shape of the potential.

The potential must be rather flat over some range, with a slight downward slope, and then at a critical value the bottom drops out. It does not fall to exactly zero. It falls to something very small. In Planck units the present vacuum energy is of order
\begin{equation}
10^{-123}.
\end{equation}
Up on the plateau it was much larger than this, but still far below the Planck scale.

\lecturetimestamp{00:43:23}
\begin{figure}[t]
\centering
\includegraphics[width=0.78\textwidth]{lecture_08_figure_03.png}
\caption{Qualitative inflaton potential with a broad plateau, a steep drop, and a low minimum.}
\end{figure}

This plot does not describe a place in space. It describes a value of the field. To say that the universe started high on the plateau is only to say that the field value, over a very large region, lay up there. It is not important if \(\phi\) varied somewhat from place to place at first, because inflation stretches those variations enormously. After enough expansion, any local patch looks almost uniform. \editorialnote{The axis labels are not legible in the blackboard frame, but the spoken discussion makes clear that the vertical axis is the potential energy \(V(\phi)\) and the horizontal axis is the field value \(\phi\).}

As long as the field sits on the broad, slowly sloping part of the potential, the potential energy behaves like vacuum energy. If \(\phi\) were perfectly frozen, it would behave exactly like vacuum energy. Leaving the omitted numerical factors aside, the relation stated in the lecture is simply
\begin{equation}
H\propto \sqrt{V(\phi)},
\end{equation}
and that gives the nearly exponential law
\begin{equation}
a(t)\propto e^{Ht}.
\end{equation}
Almost everything that can now be observed happens after the field goes over the edge. The drop is where the large vacuum-like energy on the plateau is converted into radiation, ordinary matter, dark matter, and the tiny leftover vacuum energy seen today.

\begin{classroomqa}
\lecturetimestamp{00:44:07}
\audiencequestion{What is \(\phi\) a field on, and are Einstein's equations still being used?}
\lecturerresponse{Yes. \(\phi\) is a field on spacetime, and the same Friedmann equation is still in use. In this regime the energy density on the right-hand side is supplied mainly by \(V(\phi)\), provided the field changes slowly enough.}
\end{classroomqa}

\begin{classroomqa}
\lecturetimestamp{00:47:04}
\audiencequestion{How can vacuum-like inflaton energy turn into matter and radiation?}
\lecturerresponse{When the field rolls to the bottom and oscillates, those oscillations are the quanta of the field. If the quanta are unstable, they can decay into photons, electrons, positrons, and other familiar particles.}
\end{classroomqa}

\section{Slow Roll and Hubble Friction}

Once the field goes over the edge, why should it not simply slosh forever? Why should it not swing back and forth and perhaps even climb again? The answer given here is a cosmological damping mechanism. It is not genuine microscopic friction, but it behaves that way in the equation of motion.

Ignore spatial gradients and treat the field as homogeneous. Then the kinetic term and the potential term look like a single degree of freedom with \(T-V\) dynamics, but because they are energy densities they must be multiplied by the physical volume \(a^3\). One is led to a homogeneous effective Lagrangian of the form
\begin{equation}
L_{\mathrm{hom}}=a^3(t)\left[\frac12\dot\phi^2-V(\phi)\right].
\end{equation}
\editorialnote{The top line on the board is partially cropped. This homogeneous Lagrangian is reconstructed only to the extent needed to reproduce the secure equation of motion written below it.}

The equation that is secure on the board is
\begin{equation}
\frac{d}{dt}\left(a^3\dot\phi\right)=-a^3\frac{\partial V}{\partial\phi}.
\end{equation}
Expanding the time derivative gives
\begin{equation}
a^3\ddot\phi+3a^2\dot a\,\dot\phi=-a^3\frac{\partial V}{\partial\phi},
\end{equation}
and division by \(a^3\) yields
\begin{equation}
\ddot\phi+3H\dot\phi=-\frac{\partial V}{\partial\phi}.
\end{equation}
The minus sign is corrected explicitly in the spoken derivation: the force is minus the derivative of the potential energy.

\lecturetimestamp{00:57:10}
\begin{figure}[t]
\centering
\includegraphics[width=0.78\textwidth]{lecture_08_figure_05.png}
\caption{Homogeneous inflaton equation of motion showing the Hubble-friction term \(3H\dot\phi\).}
\end{figure}

The term \(3H\dot\phi\) is what is called Hubble friction. It comes entirely from differentiating the factor \(a^3\). Dynamically, however, it looks like the drag term in the motion of a rock falling through water. There is a force term, an acceleration term, and a term linear in the velocity. The field quickly reaches a terminal velocity and then moves slowly along the plateau. That is the slow-roll regime. While the slope is shallow and the field moves slowly, the potential energy changes only slowly, so the expansion stays close to exponential.

Near the minimum the field oscillates. Those oscillations are the inflaton quanta. If the quanta are unstable, they decay, and that is how the energy on the plateau is turned into the hot matter and radiation content of the later universe.

\begin{classroomqa}
\lecturetimestamp{00:50:48}
\audiencequestion{Why doesn't the field just oscillate forever, and what provides the damping?}
\lecturerresponse{Expansion itself provides a friction-like term. The factor \(a^3\) in the homogeneous Lagrangian produces the \(3H\dot\phi\) term, and that behaves like viscous drag.}
\end{classroomqa}

\section{Quantum Fluctuations and Horizon Exit}

There is now an apparent problem. If inflation makes the universe extraordinarily smooth, why does it not erase the very seeds that galaxies need? The answer is quantum mechanics. The surviving fluctuations are not remnants of an imperfect ironing-out. They are generated anew.

Forget expansion for a moment. Even if a field were almost perfectly homogeneous, quantum mechanics would not let both \(\phi\) and \(\dot\phi\) be fixed exactly. The ground state is still agitated. This is the same reason an ordinary harmonic oscillator has zero-point energy. In everyday circumstances those fluctuations are hard to see because they are short-wavelength and rapidly oscillating, so observations average over them.

Expansion changes the story. In comoving coordinates the field mode keeps roughly the same coordinate wavelength, while the physical wavelength grows with the scale factor:
\begin{equation}
\lambda_{\rm phys}=a(t)\,\Delta x.
\end{equation}
During inflation the physical Hubble scale is approximately fixed,
\begin{equation}
d_H\sim H^{-1},
\end{equation}
so the corresponding comoving horizon shrinks like
\begin{equation}
d_{H,\mathrm{coord}}\sim \frac{1}{aH}.
\end{equation}

The rope analogy makes the freezing intuitive. A wave oscillates because each bit of rope is pulled by neighboring bits. But if the wavelength becomes as large as the horizon, the relevant neighbors are beyond causal contact. The restoring pull disappears. Once a mode is longer than the horizon, it freezes.

This happens again and again. Longer wavelengths freeze first. Shorter wavelengths freeze later. Even a mode that begins with a tiny wavelength is eventually stretched until it crosses the horizon. Thus inflation does not simply erase structure. It keeps producing frozen field variations.

The process also does not stop just because inflation lasts a long time. Old frozen modes are stretched to enormous scales, but new modes are always freezing out. The accumulated pattern becomes statistically steady in the sense relevant for cosmology. That is the origin of the approximately scale-invariant spectrum: the same horizon-crossing process repeats on successive logarithmic scales and produces roughly the same amplitude from one wavelength band to the next.

After the field rolls off the plateau, the horizon grows again. Modes that had been outside the temporary inflationary horizon later come back in. In the usual language they leave the horizon during inflation and re-enter it afterward. By then they are stretched to macroscopic scales, so when they begin to oscillate again they do so as very large structures.

\begin{classroomqa}
\lecturetimestamp{01:12:02}
\audiencequestion{What does the horizon look like on the spacetime graph?}
\lecturerresponse{The physical distance to the horizon stays roughly fixed at \(c/H\) during exponential expansion, but on a comoving graph that fixed physical distance looks smaller and smaller because the coordinate scale is being stretched. In comoving terms the horizon shrinks like \(1/(aH)\).}
\end{classroomqa}

\begin{classroomqa}
\lecturetimestamp{01:15:59}
\audiencequestion{What happens to the amplitude as the fluctuations are stretched, and is there a steady state?}
\lecturerresponse{Individual modes dilute as they are stretched, but new modes keep freezing out. The result is a statistical steady pattern, approximately scale invariant, rather than a complete disappearance of fluctuations.}
\end{classroomqa}

\section{From Frozen Field Variations to Density Variations}

Once the field has acquired small variations from place to place, those variations are converted into density differences when different regions go over the edge at slightly different times. A region a little farther downhill reaches the cliff first. If its vacuum energy is converted into radiation, that radiation dilutes with expansion. A neighboring region still on the plateau remains vacuum dominated and does not dilute. By the time the second region falls, the first one has already redshifted part of its energy away.

That is why slower roll gives larger final density contrasts. If the field moves only very slowly along the plateau, then a small difference in field value corresponds to a larger time lag between one region and the next reaching the cliff. The earlier-falling region has more time to dilute before the later-falling region converts its energy. The flatter the slope, the smaller the terminal velocity, and the larger the final density variation. A higher plateau also tends to give larger fluctuations. The width of the plateau mainly controls how long inflation lasts.

Only a combination of height and slope is tightly constrained by the fluctuation data. If the energy scale is high enough, there should also be a visible gravitational-wave contribution. The numerical bounds are given only as rough ballpark estimates. Roughly speaking, if the energy density on the plateau were much above about \(10^{-14}\) in Planck units, the microwave background would already show too much gravitational-wave structure; if it were much below about \(10^{-16}\), there would be little realistic chance of seeing such a signal. \editorialnote{The speaker explicitly treats these numbers as approximate and says he does not remember the precise values. They are kept here only as order-of-magnitude statements.}

The units are explained in the usual natural-units way. The Compton wavelength of a particle of mass \(m\) is
\begin{equation}
\lambda_C=\frac{\hbar}{mc}.
\end{equation}
If \(c=\hbar=1\), then mass and inverse length have the same units. Energy density is energy per volume, so it has units of mass to the fourth:
\begin{equation}
[\rho]=(\mathrm{mass})^4.
\end{equation}

Seen at a single point, the accumulation of frozen modes looks like a random walk. Different points in space undergo different random walks, and that is precisely what turns the quantum noise into spatial variation.

\begin{classroomqa}
\lecturetimestamp{01:33:04}
\audiencequestion{Does \(\phi\) undergo a random walk, and does that make the earlier starting point irrelevant?}
\lecturerresponse{At any fixed point the successive frozen modes add with random phase, so \(\phi\) does perform a random walk. That very success is also the difficulty: inflation keeps replacing earlier information with later fluctuations, so the stage before inflation becomes hard to recover.}
\end{classroomqa}

\section{What Inflation Hides}

The theory is successful in a very specific sense. It correlates fluctuations over a huge range of scales with one underlying mechanism. But the same mechanism washes out what came before. Earlier irregularities are stretched away, and the quantum fluctuations produced later are the ones that remain relevant. That is why inflation is both attractive and frustrating. It explains the smoothness and the fluctuation spectrum, but it also erases much of the earlier history.

This leads to an uneasiness about the starting point. A very high energy density is not the problem. The problem is that the plateau energy seems neither naturally Planckian nor naturally tiny. It is high in ordinary units, but still small in Planck units. It does not look like an obvious first stopping place for a fundamental theory.

There is also no literal \(t=0\) in this simple picture. Once inflation has lasted too long, almost all information about any earlier stage is gone. If the number of e-foldings were only barely above the minimum needed, perhaps some trace of the previous stage could survive, but in general inflation wipes out the past.

The same caution appears in the treatment of de Sitter space. The pure exponential law is correct for the inflationary branch at late times, but the exact de Sitter solution is better described by
\begin{equation}
a(t)\propto \cosh(Ht).
\end{equation}
That solution contracts to a minimum size and then expands. Whether the contracting branch should be taken literally in cosmology is left open.

\begin{classroomqa}
\lecturetimestamp{01:37:31}
\audiencequestion{Could there have been no start, with exponential behavior continuing indefinitely into the past?}
\lecturerresponse{Not in this simple picture. Exact de Sitter space is not an eternally expanding exponential for all time; it is more like \(\cosh(Ht)\), with a contracting branch and an expanding branch. Whether that earlier branch should be taken literally is left as an open question.}
\end{classroomqa}

\section{Tuning, Curvature, and Anthropic Bounds}

To first approximation the important parameters are the height of the plateau, its slope, and its width \(\Delta\phi\). The slope matters because it sets the terminal velocity and therefore the size of the fluctuations. The width matters because it helps determine how long inflation lasts. A common prejudice is that a basic field excursion should not be much larger than the Planck mass, and that the energy density should not exceed order one in Planck units. Under such rough bounds, enough inflation does not fill the whole parameter space. The estimate quoted here is a few-percent tuning, about \(2\%\), to make the slope shallow enough for roughly \(60\) to \(62\) e-foldings.

\begin{classroomqa}
\lecturetimestamp{01:46:02}
\audiencequestion{Can inflation be understood in terms of a parameter space, and how tuned is it?}
\lecturerresponse{Yes. To first approximation one thinks in terms of the height, slope, and width of the plateau. The lecture gives a rough estimate that getting enough e-foldings requires tuning at the level of a few percent, about \(2\%\).}
\end{classroomqa}

That mild tuning is then compared with the much sharper puzzle of the present cosmological constant, of order \(10^{-123}\). The discussion turns to anthropic questions because both problems can be reframed by asking what is required for structure to form.

A positive cosmological constant acts, in the Newtonian picture, like a universal repulsion proportional to distance. If it were much larger, it would not have to blow galaxies apart today in order to be fatal. It would only have to be large enough to stop the primordial \(10^{-5}\) fluctuations from ever condensing into galaxies. The stated example is that a value roughly two orders of magnitude above the observed one would already do that. This is the point associated with Weinberg's argument: the cosmological constant need not be exactly zero; it need only be small enough to allow structure.

Negative curvature works in the same direction. In the Newtonian language of the earlier lectures, \(k<0\) means recession faster than escape velocity. Too much outward thrust also prevents the small primordial overdensities from ever turning into galaxies. Inflation dilutes that curvature. The quoted lower bound is about \(59.5\) e-foldings. A value around \(62\) is then taken as a representative inflationary history safely above the threshold.

This changes the probability question. If one asks for the raw volume of parameter space that gives enough inflation, the answer may be only a few percent. If one conditions on being in a universe where galaxies form and observers exist to ask the question, the conditional probability can become much larger. That is where the discussion closes: not with a final solution, but with the fine-tuning puzzles, the anthropic reinterpretation of them, and the recognition that this is where the great debates remain.
